\documentclass[11pt,letterpaper]{article}
\usepackage[T1]{fontenc}
\usepackage{lmodern}
\usepackage[margin=0.95in,headheight=15pt]{geometry}
\usepackage{amsmath,amssymb,amsthm,mathtools,booktabs,array,longtable}
\usepackage{microtype,enumitem,xurl,hyperref,fancyhdr}
\hypersetup{hidelinks,pdftitle={Contact-Complete Multishift Lerch Reflection},pdfauthor={Research continuation prepared for Vladimir Reshetnikov}}
\pagestyle{fancy}\fancyhf{}
\fancyhead[L]{\small Multishift Lerch reflection}
\fancyhead[R]{\small ProveIt research continuation}
\fancyfoot[C]{\thepage}
\setlength{\parskip}{3pt}\setlength{\parindent}{15pt}
\setlist{itemsep=3pt,topsep=5pt}
\allowdisplaybreaks[2]\numberwithin{equation}{section}
\newtheorem{theorem}{Theorem}[section]
\newtheorem{lemma}[theorem]{Lemma}
\newtheorem{proposition}[theorem]{Proposition}
\newtheorem{corollary}[theorem]{Corollary}
\theoremstyle{definition}\newtheorem{definition}[theorem]{Definition}
\theoremstyle{remark}\newtheorem{remark}[theorem]{Remark}
\newcommand{\C}{\mathbb C}\newcommand{\R}{\mathbb R}\newcommand{\N}{\mathbb N}
\newcommand{\Rea}{\operatorname{Re}}\newcommand{\Ima}{\operatorname{Im}}
\newcommand{\Li}{\operatorname{Li}}\newcommand{\Fp}{\operatorname{Fp}}
\newcommand{\Res}{\operatorname{Res}}\newcommand{\Log}{\operatorname{Log}}
\newcommand{\dd}{\,\mathrm d}\newcommand{\M}{\mathcal M}
\newcommand{\F}{\mathcal F}\newcommand{\B}{\mathcal B}
\newcommand{\stir}[2]{\genfrac{[}{]}{0pt}{}{#1}{#2}}
\newcommand{\src}[1]{\nolinkurl{#1}}
\newcommand{\pin}{\src{f7e35be357654198c1546fd40f363694c5bcc2cc}}
\title{\textbf{Contact-Complete Multishift Lerch Reflection}\\[6pt]
\large Independent spectral jets, generalized Stieltjes log moments,\\
confluent contacts, and Gamma--polylogarithm identities}
\author{Research continuation prepared for Vladimir Reshetnikov's ProveIt programme\\
\small Mathematical derivations and computational checks with OpenAI assistance}
\date{10 October 2026}
\begin{document}
\maketitle
\begin{abstract}
We develop an exact centered-contour completion for the multiplier
$(\pi\csc\pi p)^R\prod_j(p+a_j-1)^{-s_j}$, with independent complex
shifts of positive real part, independent entire spectral orders, and
$\Rea R>0$. Horizontal cuts give a finite, explicit collection of correction
integrals. When cuts have the same height their jumps combine, rather than
being added as isolated one-factor corrections. This resolves the
independent-shift centered-reflection problem posed in the incoming
\emph{Continuous-Depth Lerch Calculus} report, in the precise contour and
branch conventions stated here. Collisions have a unique holomorphic
completion, including arbitrary simultaneous collisions and centered poles.

At zero spectral orders every mixed derivative becomes an ordinary
logarithmic integral with an explicit finite jump polynomial. The cubic
polynomial contains a non-additive $-\pi^2$ overlap term. A normally
convergent Stirling compiler expresses the resulting depth-two moments
using one generalized Stieltjes constant and convergent Hurwitz-zeta jets.
Balanced log ratios give exact harmonic-number generating families.
At positive integer orders, Hermite residues and half-residue contacts
replace branch-cut integrals; all spectral derivatives of the local
contact distribution are given explicitly.

A second continuation supplies every order derivative on the odd and even
centered evaluation lattices. Finite Bell formulas retain the necessary
generalized alternating-Hurwitz jets and therefore do not infer derivatives
from discrete values. Consequences include entire-order Gamma-weighted
polylogarithm identities, explicit Stieltjes-valued order derivatives,
polygamma weights, and anchored primitives through every integration
resonance. The classical Gamma transform, one-shift hyperbolic Stieltjes
formula, and existing single-anchor expansion are credited as inputs.
No global priority claim, arithmetic independence assertion, or proof of
the separate Gaussian $S_6$ and $S_8$ candidates is made.
\end{abstract}
\vspace{4pt}
\noindent\textbf{Repository anchor.} \pin. The repository was read only.
This package is an additive research proposal, not a rebuilt canonical manuscript.
\clearpage
\begingroup\small\setlength{\parskip}{0pt}\tableofcontents\endgroup
\clearpage

\section{Research target and contribution boundaries}
\label{mr:sec:scope}
The canonical ProveIt manuscript and its incoming queue contain a substantial
body of exact polylogarithm, harmonic-sum, Gamma, and Stieltjes calculus
\cite{Repo,CDL,Tilt}. The aim here is not to repeat that infrastructure.
Two questions in the incoming \emph{Continuous-Depth Lerch Calculus}
report \cite[Section 13, Q4 and Q8]{CDL} are particularly concrete:
independent-shift centered reflection with collision and branch rules;
and genuine spectral derivatives on the centered odd evaluation lattice.
The former report's section numbers refer to its delivered source; the
question labels are the stable identifiers used in our claim ledger.

We answer the first by a horizontal-cut formula valid for arbitrary complex
shifts in the right half-plane. It remains an identity between the stated
contour transform, a centered integral, and finitely many explicit cut
integrals: it is \emph{not} a claim of finite reduction to ordinary Lerch
functions at every continuous depth. We answer the second by explicit
finite Bell formulas in derivatives of the already defined entire family
$Z_R$. This is an exact representation of every jet, not an evaluation
of each retained $Z_R$ derivative in a smaller arithmetic constant space.

The Gamma-product Fourier transform is classical \cite{CS,Beta}.
The one-shift hyperbolic formula for generalized Stieltjes constants, with
$\Rea a>1/2$, is explicitly present in Johansson--Blagouchine
\cite[equation (2)]{JB}. The independent-shift anchor expansion is already
part of the incoming tilted-Lerch work \cite{Tilt}. Residues, Hermite
divided differences, the Plemelj boundary rule, and Taylor subtraction are
standard mechanisms. Our additions are their complete joint use in the
specified multishift family, the all-index mixed jump and Stieltjes
compiler, and the explicit resonant spectral jets and consequences.
Historical priority for each specialization has not been established.

\subsection{Conventions}
We use the principal logarithm off the nonpositive real axis and its
upper/lower limiting values on that axis. Set
\begin{equation}\label{mr:eq:g}
 g(p)=\frac\pi{\sin\pi p},\qquad 0<\Rea p<1.
\end{equation}
Here $\Rea g(p)>0$, so $g(p)^R=\exp(R\Log g(p))$ has a fixed
holomorphic logarithm on the strip, real on $(0,1)$.
Spectral derivatives are denoted by $f^{[q]}(s)=\partial_s^q f(s)$;
ordinary primes on $\psi$ denote argument derivatives. We normalize
\begin{equation}\label{mr:eq:stieltjes}
 \zeta(1+w,a)=\frac1w+\sum_{m\ge0}\frac{(-1)^m\gamma_m(a)}{m!}w^m,
 \qquad \gamma_0(a)=-\psi(a),\qquad\Rea a>0.
\end{equation}
The ordinary Lerch convention is
$\Phi(z,s,a)=\sum_{n\ge0}z^n(n+a)^{-s}$ in its defining region,
continued with the right-half-plane logarithms of $n+a$ \cite{Lerch,Zeta}.
The exponential Bell polynomials satisfy
\begin{equation}\label{mr:eq:bell}
 \exp\left(\sum_{k\ge1}d_k\frac{z^k}{k!}\right)
 =\sum_{q\ge0}\B_q(d_1,\ldots,d_q)\frac{z^q}{q!},\qquad\B_0=1.
\end{equation}
We write $H_n^{(k)}=\sum_{j=1}^n j^{-k}$, $H_0^{(k)}=0$,
and $\stir{k}{r}$ for unsigned Stirling numbers of the first kind,
including $\stir00=1$ and $\stir{k}{0}=0$ for $k>0$.

\section{The multishift entire-order transform}
\label{mr:sec:master}
For $r\ge1$, $\Rea R>0$, and $\Rea a_j>0$, define
\begin{equation}\label{mr:eq:master}
 \M_R(\boldsymbol s;\boldsymbol a)=\frac1{2\pi i}
 \int_{c-i\infty}^{c+i\infty}g(p)^R
       \prod_{j=1}^r(p+a_j-1)^{-s_j}\dd p,
 \qquad\max(0,1-\min_j\Rea a_j)<c<1.
\end{equation}
Every resolvent on this line has positive real part. Repeated shifts
may be merged by adding their orders; this identity is literal, not a
choice of continuation after collision.

\begin{theorem}[Holomorphy and ordinary convolution]\label{mr:thm:master}
The transform \eqref{mr:eq:master} is independent of its allowed line,
entire in all $s_j$, and jointly holomorphic for $\Rea R>0$ and
$\Rea a_j>0$. Every fixed mixed parameter derivative can be taken
under the integral. If $R=\rho_1+\cdots+\rho_r$ with $\Rea\rho_j>0$,
then it is the ordinary absolutely convergent convolution at zero of
$\F_{\rho_j,a_j}(s_j;\cdot)$, where
\begin{equation}\label{mr:eq:F}
 \F_{\rho,a}(s;x)=\frac1{2\pi i}\int_{c-i\infty}^{c+i\infty}
 e^{px}g(p)^\rho(p+a-1)^{-s}\dd p.
\end{equation}
All fixed mixed parameter derivatives of this convolution converge
absolutely as well.
\end{theorem}
\begin{proof}
On any allowed line and compact parameter set, $\Log g(c+it)$ has
bounded imaginary part and real part $-\pi|t|+O(1)$.
Consequently the integrand, after any fixed number of parameter or
$t$ derivatives, is bounded by a polynomial in $1+|t|$ times
$e^{-\pi\Rea R|t|}$. Spectral differentiation inserts logarithms,
shift differentiation inserts inverse resolvents, and depth differentiation
inserts $\Log g$, all covered by this estimate. Uniform convergence
proves joint holomorphy. Cauchy's theorem on rectangles proves line
independence, because the horizontal sides tend to zero exponentially.

For a common allowed $c$, the inverse Fourier transform of
$g(c+it)^{\rho_j}(c+it+a_j-1)^{-s_j}$ is
$e^{-cx}\F_{\rho_j,a_j}(s_j;x)$, a Schwartz function. Fourier
multiplication proves the convolution identity. On the hyperplane
$x_1+\cdots+x_r=0$ the exponential weights cancel. Absolute integration
is bounded by $r-1$ weighted $L^1$ norms and one weighted $L^\infty$
norm. The same estimates are locally uniform after every fixed mixed
derivative, giving the asserted differentiation rule.
\end{proof}

For reference, the classical beta transform gives
\begin{align}
 K_R(x)&=\frac{(2\pi)^{R-2}e^{x/2}}{\Gamma(R)}G_R(x),
 &G_R(x)&=\Gamma\left(\frac R2+\frac{ix}{2\pi}\right)
          \Gamma\left(\frac R2-\frac{ix}{2\pi}\right),\label{mr:eq:kernel}\\
 b_R(x)&=e^{-x/2}K_R(x),
 &\int_\R e^{-px}K_R(x)\dd x&=g(p)^R.\label{mr:eq:transform}
\end{align}
For complex $R$, $G_R$ is a product, not a squared modulus.
Indeed substitution $t=e^u$ in Euler's beta integral proves
\begin{equation}\label{mr:eq:beta}
 \frac{\Gamma(R/2+iv)\Gamma(R/2-iv)}{\Gamma(R)}
 =\int_\R e^{ivu}(2\cosh(u/2))^{-R}\dd u.
\end{equation}
Fourier inversion and Gamma asymptotics give \eqref{mr:eq:transform};
these are the classical inputs, also used in \cite{CDL,CS}.
When $\Rea s_j>0$, Laplace integration yields the useful ordinary formula
\begin{equation}\label{mr:eq:fractional}
 \M_R(\boldsymbol s;\boldsymbol a)=
 \frac1{\prod_j\Gamma(s_j)}\int_{(0,\infty)^r}
 b_R(U)\prod_j u_j^{s_j-1}e^{-(a_j-1/2)u_j}\dd\boldsymbol u,
 \qquad U=\sum_j u_j.
\end{equation}
The large-$U$ exponential is controlled by $\min\Rea a_j>0$.
This formula must not be used as an improper integral at negative orders.

\begin{lemma}[Depth-two scalar normalization]\label{mr:lem:E}
Writing $E_a(s)=\M_2(s;a)$, one has, as an entire identity,
\begin{equation}\label{mr:eq:E}
 E_a(s)=s\zeta(s+1,a),\qquad E_a(0)=1,
 \qquad E_a^{[q]}(0)=q(-1)^{q-1}\gamma_{q-1}(a)\quad(q\ge1).
\end{equation}
\end{lemma}
\begin{proof}
Gamma reflection gives $K_2(-u)=u/(e^u-1)$.
The one-variable fractional formula is therefore
$\Gamma(s)^{-1}\int_0^\infty u^s e^{-au}(1-e^{-u})^{-1}\dd u$.
Expanding geometrically for $\Rea s>0$ gives $s\zeta(s+1,a)$.
Entire continuation follows from Theorem~\ref{mr:thm:master}; the
Taylor coefficients follow from \eqref{mr:eq:stieltjes}.
\end{proof}

\section{Horizontal cuts and complete centered reflection}
\label{mr:sec:cuts}
Put $p_j=1-a_j=x_j+iy_j$. Temporarily assume $x_j\ne1/2$.
The central integral is
\begin{equation}\label{mr:eq:J}
 J_R(\boldsymbol s;\boldsymbol a)=\frac1{2\pi}\int_\R
 \left(\frac\pi{\cosh\pi t}\right)^R
 \prod_j(a_j-1/2+it)^{-s_j}\dd t.
\end{equation}
It is split at $t=y_j$ when a principal cut meets the line. No single
point value at a jump affects this integral. It is entire in
$\boldsymbol s$ for each fixed configuration with $x_j\ne1/2$.

For each distinct height $y$ let $I_y=\{j:y_j=y\}$ and
$X_y=\max_{j\in I_y}x_j$. Only heights with $X_y>1/2$ contribute.
For $1/2<x<X_y$, off the finitely many $x_j$, set
\begin{equation}\label{mr:eq:active}
 A_y(x)=\{j\in I_y:x<x_j\},\quad
 \sigma_y(x)=\sum_{j\in A_y(x)}s_j,
\end{equation}
\begin{equation}\label{mr:eq:D}
 D_y(x)=\prod_{j\in I_y}|x-x_j|^{-s_j}
         \prod_{j\notin I_y}(x+iy-p_j)^{-s_j}.
\end{equation}
The latter factors use their principal logs; their imaginary parts are
nonzero. Define, splitting at all $x_j$ on the integration segment,
\begin{equation}\label{mr:eq:C}
 C_R(\boldsymbol s;\boldsymbol a)=\frac1\pi
 \sum_{y:X_y>1/2}\int_{1/2}^{X_y}
 g(x+iy)^R\sin(\pi\sigma_y(x))D_y(x)\dd x.
\end{equation}

\begin{theorem}[Complete multishift centered reflection]\label{mr:thm:cut}
Suppose $\Rea R>0$, $\Rea a_j>0$, and no $x_j=1/2$.
For each block $B$ of identical branch points to the right of the
central line, initially require $\Rea\sum_{j\in B}s_j<1$.
Then all integrals below are ordinary and absolutely convergent, and
\begin{equation}\label{mr:eq:complete}
 \boxed{\M_R(\boldsymbol s;\boldsymbol a)
       =J_R(\boldsymbol s;\boldsymbol a)
        +C_R(\boldsymbol s;\boldsymbol a).}
\end{equation}
The \emph{sum} of cut terms has a unique entire continuation in all
spectral orders. The completed expression is jointly holomorphic in
all right-half-plane shifts, including every shift collision. Formula
\eqref{mr:eq:C}, not a sum of independently continued one-factor
corrections, specifies the branch convention when cuts overlap.
\end{theorem}
\begin{proof}
Choose the original line to the right of both $1/2$ and every $x_j$.
In the intervening strip, remove the horizontal cut segments ending
at the $p_j$, and integrate around a truncated rectangle with thin
slits. On a segment at height $y$, the boundary values of an active
resolvent are $|x-x_j|^{-s_j}e^{\mp i\pi s_j}$.
Inactive same-height factors and all other-height factors have no jump
across that segment. The lower value minus the upper value of the
product is therefore $2i\sin(\pi\sigma_y(x))D_y(x)$.
The original upward line minus the centered upward line equals the
integral of this lower-minus-upper jump from left to right, divided by
$2\pi i$. This proves the sign and normalization in \eqref{mr:eq:C}.

Small circles around a block $B$ have size
$O(\epsilon^{1-\Rea\sum_{B}s_j})$, and hence vanish in the stated
initial domain. Away from these points the slit limits are bounded.
The horizontal sides at large height vanish by the exponential estimate
used in Theorem~\ref{mr:thm:master}. These facts justify every limiting
operation in the contour deformation.

For fixed shifts, both $\M_R$ and $J_R$ are entire in the orders;
thus their difference gives the unique entire continuation of the
\emph{complete} cut sum. A constructive continuation subtracts local
Taylor coefficients of the analytic factor at each endpoint. The
subtracted terms are of the form
$h_k\epsilon^{k+1-\sigma}/(k+1-\sigma)$ and the remaining integral
converges in successively larger order domains. Terms at a common
endpoint must be combined before evaluating a pole. Uniqueness follows
from agreement in an open polydomain.

Finally, $\M_R$ itself is jointly holomorphic in the shifts, including
collisions, by Theorem~\ref{mr:thm:master}. This is a genuine analytic
extension of the completed expression, not an assertion that an
individual cut term or $J_R$ is holomorphic under every change of cut
geometry. It fixes every limiting value uniquely.
\end{proof}

\subsection{Why colliding cut heights do not create ambiguity}
The requisite algebra is the telescoping product identity
\begin{equation}\label{mr:eq:telescope}
 \prod_{j=1}^r u_j-\prod_{j=1}^r v_j
 =\sum_{j=1}^r (u_j-v_j)
       \prod_{k<j}u_k\prod_{k>j}v_k.
\end{equation}
As distinct horizontal heights approach one another, the other factors
on a cut acquire the appropriate upper or lower limiting phases.
Their sum telescopes to the single jump of the full product. Any ordering
of the approaching heights gives the same merged jump. Equality of
branch points then replaces the relevant orders by their sum.
No assumption on ratios of distances in a simultaneous collision is
needed for the completed identity. The contour at the right remains
uniformly separated from every branch point on a compact parameter set.

\subsection{The real-shift specialization}
For real $a_j>0$, there is one height, $y=0$. Thus
\begin{equation}\label{mr:eq:realcut}
 C_R=\frac1\pi\int_{1/2}^{\max_j(1-a_j)}
 g(x)^R\sin\left(\pi\sum_{x<1-a_j}s_j\right)
 \prod_j|x+a_j-1|^{-s_j}\dd x,
\end{equation}
with the integral zero if its upper endpoint is at most $1/2$.
In particular, no correction is present when every $a_j>1/2$.
At $a_j=1/2$, the central power may cease to be locally integrable.
The holomorphic value is still defined by \eqref{mr:eq:master}; its
integer-order boundary convention is given in Section~\ref{mr:sec:contact}.
The zero-order logarithmic jets, treated next, need no finite parts even
at centered shifts.

\section{All mixed logarithmic moments and overlap corrections}
\label{mr:sec:logs}
Let $\boldsymbol q=(q_1,\ldots,q_r)$ be nonnegative integers and
$Q=\sum_jq_j$. On a cut at height $y$ define
\begin{equation}\label{mr:eq:Lpm}
 L_{j,y}^{\pm}(x)=
 \begin{cases}
 \log|x-x_j|\pm i\pi,&j\in A_y(x),\\
 \Log(x+iy-p_j),&j\notin A_y(x),
 \end{cases}
\end{equation}
where the second line is unambiguous off the cut endpoints: it is a
real log for inactive factors at the same height. Set
\begin{equation}\label{mr:eq:jumpP}
 P_{\boldsymbol q,y}(x)=\frac1{2i}
 \left\{\prod_jL_{j,y}^+(x)^{q_j}
       -\prod_jL_{j,y}^-(x)^{q_j}\right\}.
\end{equation}
For complex shifts or depth this is generally complex; it must not be
replaced by the imaginary part of the first product.
Define the completed logarithmic moment
\begin{align}
 \Lambda_{R,\boldsymbol q}(\boldsymbol a)
 ={}&\frac1{2\pi}\int_\R
 \left(\frac\pi{\cosh\pi t}\right)^R
 \prod_j\Log(a_j-1/2+it)^{q_j}\dd t\nonumber\\
 &-\frac1\pi\sum_{y:X_y>1/2}\int_{1/2}^{X_y}
       g(x+iy)^R P_{\boldsymbol q,y}(x)\dd x.
 \label{mr:eq:Lambda}
\end{align}
As usual a factor with $q_j=0$ is one, even at its isolated endpoint.

\begin{theorem}[Ordinary all-index mixed jets]\label{mr:thm:logs}
For $\Rea R>0$ and arbitrary $\Rea a_j>0$, including centered and
coincident shifts, every integral in \eqref{mr:eq:Lambda} is ordinary
and absolutely convergent. The identity
\begin{equation}\label{mr:eq:logjet}
 \boxed{\Lambda_{R,\boldsymbol q}(\boldsymbol a)
 =(-1)^Q\left.\partial_{s_1}^{q_1}\cdots\partial_{s_r}^{q_r}
             \M_R(\boldsymbol s;\boldsymbol a)\right|_{\boldsymbol s=0}}
\end{equation}
holds. In particular the completed moment is holomorphic in all the
shifts, even though its separate raw and cut terms need not be.
\end{theorem}
\begin{proof}
First take noncentered shifts and orders in a sufficiently small
polydisk about zero. Theorem~\ref{mr:thm:cut} is then an ordinary
identity, and order derivatives are dominated near each endpoint by
$u^{-\epsilon}(1+|\log u|)^Q$ with some $\epsilon<1$.
The original and centered contours have exponential decay at infinity.
Thus differentiation is justified termwise. The upper/lower product
rule in \eqref{mr:eq:jumpP} is exactly the derivative of the jump of the
powers. There are $Q$ minus signs from spectral differentiation;
rearranging them gives the \emph{minus} sign in \eqref{mr:eq:Lambda}.

At centered or colliding shifts, only logarithmic endpoint singularities
remain. On each compact interval every fixed power of a translated
logarithm is locally integrable, uniformly under small translations;
H\"older's inequality gives the corresponding product estimate.
One way to see uniform integrability explicitly is to choose any
$\epsilon>0$ and dominate $|\log|u-h||^q$ near its singularity by
$C_{q,\epsilon}|u-h|^{-\epsilon}$, choosing the exponent after the
finite number of factors is fixed. Away from small endpoint
neighborhoods the boundary values converge uniformly. When cut
heights merge, \eqref{mr:eq:telescope} combines their jumps. These facts
pass the ordinary log formula to all the stated collisions. Equivalently,
the right contour, uniformly separated from the shifts, gives their
unique common limit. Equality with the holomorphic right side then
proves the last assertion.
\end{proof}

\subsection{An explicit finite polynomial for real shifts}
For real shifts put $\ell_j=\log|x-p_j|$ and
$\chi_j=\mathbf1_{x<p_j}$. Expansion of \eqref{mr:eq:jumpP} gives
\begin{equation}\label{mr:eq:jumpcompiler}
 P_{\boldsymbol q}(x)=
 \sum_{\substack{0\le k_j\le q_j\\ K=\sum_jk_j\text{ odd}}}
 (-1)^{(K-1)/2}\pi^K
 \prod_j\binom{q_j}{k_j}\ell_j^{q_j-k_j}\chi_j^{k_j}.
\end{equation}
This is a terminating exact compiler at every harmonic degree.
In the first two mixed cases it reads
\begin{align}
 \frac1\pi P_{(1,1)}&=\chi_1\ell_2+\chi_2\ell_1,\label{mr:eq:pairP}\\
 \frac1\pi P_{(1,1,1)}
 &=\chi_1\ell_2\ell_3+\chi_2\ell_1\ell_3+\chi_3\ell_1\ell_2
   -\pi^2\chi_1\chi_2\chi_3.\label{mr:eq:tripleP}
\end{align}
The last term is essential. It is present on the common portion of
three cuts, and cannot be recovered by adding independent pairwise
corrections. At higher degree, every odd intersection order contributes
as specified by \eqref{mr:eq:jumpcompiler}.

\section{Generalized Stieltjes identities on the full right half-plane}
\label{mr:sec:stieltjes}
At a common shift the orders combine, and
Lemma~\ref{mr:lem:E} immediately gives an all-index consequence.
\begin{corollary}[Diagonal completed moments]\label{mr:cor:diag}
For $Q\ge1$ and $\Rea a>0$,
\begin{equation}\label{mr:eq:diagonal}
 \Lambda_{2,\boldsymbol q}(a,\ldots,a)=-Q\gamma_{Q-1}(a).
\end{equation}
For $Q=0$ the value is one.
\end{corollary}

For real $a>0$ the formula can be displayed without multishift notation.
For every $m\ge0$,
\begin{align}
 \gamma_m(a)={}&-\frac{\pi}{2(m+1)}\int_\R
 \frac{\Log(a-1/2+it)^{m+1}}{\cosh^2\pi t}\dd t\nonumber\\
 &+\frac{\pi}{m+1}\mathbf1_{a<1/2}
 \int_0^{1/2-a}\csc^2\!\bigl(\pi(a+u)\bigr)
       \Ima(\log u+i\pi)^{m+1}\dd u.
 \label{mr:eq:stieltjesall}
\end{align}
The imaginary part in this \emph{real-shift} formula is legitimate.
It is not the convention for general complex shifts; for those,
\eqref{mr:eq:Lambda} with the difference of boundary values applies.
When $a>1/2$, the second line vanishes and the first line is the
Johansson--Blagouchine formula \cite[equation (2)]{JB}. Its value at
$a=1/2$ is the ordinary logarithmic limit. The extra line gives the
explicit crossing completion for $0<a<1/2$.

\subsection{A diagnostic digamma reflection}
At $m=0$, let
\begin{equation}\label{mr:eq:rawlog}
 A(a)=\frac\pi2\int_\R\frac{\Log(a-1/2+it)}{\cosh^2\pi t}\dd t.
\end{equation}
Then
\begin{equation}\label{mr:eq:digamma}
 A(a)=\psi(a)+\mathbf1_{a<1/2}\pi\cot(\pi a)
 =\begin{cases}\psi(a),&a\ge1/2,\\\psi(1-a),&0<a<1/2.\end{cases}
\end{equation}
The first equality follows by elementary integration of the correction;
the second is the classical digamma reflection formula \cite{Polygamma}.
In particular, continuing the \emph{raw} integral as though it were
$\psi(a)$ below the center gives an error of exactly $\pi$ at $a=1/4$.
This is an explicit counterexample to an invalid extension, not an error
attributed to \cite{JB}, which states its half-plane restriction.

For the cubic diagonal moment, $0<a<1/2$ gives
\begin{equation}\label{mr:eq:cubicdiag}
 \frac\pi2\int_\R\frac{\Log(a-1/2+it)^3}{\cosh^2\pi t}\dd t
 -\int_0^{1/2-a}g(a+u)^2(3\log^2u-\pi^2)\dd u
 =-3\gamma_2(a).
\end{equation}
If the $-\pi^2$ overlap is omitted inside the second integral, the
left side changes by $-\pi^3\cot(\pi a)$. The sign is fixed, for
example, at $a=1/4$ where the change is $-\pi^3$.

\section{Collision series and a finite Stieltjes-jet compiler}
\label{mr:sec:compiler}
The incoming tilted-Lerch continuation already constructs a single-anchor
expansion for unequal shifts \cite{Tilt}. We record its specialization
and proof to make the mixed-jet consequence independent of that source's
notation. This expansion itself is not counted as a new result.

\begin{proposition}[Normally convergent anchor expansion]\label{mr:prop:anchor}
Let $a>0$, $a_j=a+h_j$, and $\max_j|h_j|<a$. Put
$S=\sum_js_j$ and $K=\sum_jk_j$. Then
\begin{equation}\label{mr:eq:anchor}
 \M_R(\boldsymbol s;a+\boldsymbol h)=
 \sum_{\boldsymbol k\ge0}\frac{(-1)^K}{\prod_jk_j!}
 \prod_j(s_j)_{k_j}h_j^{k_j}\,
 \F_{R,a}(S+K;0).
\end{equation}
The series is locally normally convergent in all orders and in the
stated shift polydisk, and all fixed parameter derivatives are allowed.
\end{proposition}
\begin{proof}
Choose $c<1$ sufficiently close to one that
$\max|h_j|<c+a-1$. On the contour,
$|h_j/(p+a-1)|\le\theta<1$. Apply the binomial series to each
$(p+a+h_j-1)^{-s_j}$. On compact order sets the binomial coefficients
are at most a fixed power of $k_j+1$; the geometric factor $\theta^{k_j}$
therefore gives absolute uniform convergence. The remaining contour
integrand is exponentially decreasing, uniformly on those sets.
Integrate termwise. Holomorphic normal convergence, or the same estimates
with logarithmic factors, justifies every fixed derivative.
\end{proof}

\begin{theorem}[All mixed Stieltjes-jet coefficients]\label{mr:thm:compiler}
Under Proposition~\ref{mr:prop:anchor}, for arbitrary
$\boldsymbol q\ge0$, $Q=\sum q_j$,
\begin{align}
 \Lambda_{R,\boldsymbol q}(a+\boldsymbol h)
 ={}&(-1)^Q\sum_{\boldsymbol k\ge0}
       \frac{(-1)^K\boldsymbol h^{\boldsymbol k}}{\boldsymbol k!}
 \sum_{0\le r_j\le q_j}
 \left\{\prod_j\frac{q_j!}{(q_j-r_j)!}\stir{k_j}{r_j}\right\}
 \F_{R,a}^{[Q-|\boldsymbol r|]}(K;0).
 \label{mr:eq:compiler}
\end{align}
For $R=2$ and $Q\ge1$ the $\boldsymbol k=0$ term is exactly
$-Q\gamma_{Q-1}(a)$. Every other term is a finite combination of
Hurwitz-zeta order derivatives at positive integers at least two:
\begin{equation}\label{mr:eq:positivejets}
 E_a^{[m]}(K)=K\zeta^{[m]}(K+1,a)
             +m\zeta^{[m-1]}(K+1,a),\qquad K\ge1,
\end{equation}
with the second term omitted at $m=0$.
\end{theorem}
\begin{proof}
The coefficient identity
$(s)_k=\sum_r\stir{k}{r}s^r$ gives
$\partial_s^r(s)_k|_0=r!\stir{k}{r}$.
In differentiating a product term of \eqref{mr:eq:anchor}, choose $r_j$
derivatives to fall on the $j$th Pochhammer factor; the remaining
$Q-|\boldsymbol r|$ derivatives all act on its single argument $S+K$.
The Leibniz factor is
$\binom{q_j}{r_j}r_j!=q_j!/(q_j-r_j)!$.
This proves \eqref{mr:eq:compiler}, using the justified differentiation
of Proposition~\ref{mr:prop:anchor}. The case $K=0$ permits only
$\boldsymbol r=0$ and is \eqref{mr:eq:E}. For $K>0$, differentiating
$s\zeta(s+1,a)$ proves \eqref{mr:eq:positivejets}.
\end{proof}

This gives exactly \emph{one} anchor Stieltjes constant in this coordinate
system. It does not assert that the remaining convergent series is a
finite arithmetic reduction, or that its entries are independent.
The polydisk includes arbitrary collisions, and is not constrained by
ratios of differences between individual shifts.

For instance, with $b=a+h$ and $|h|<a$,
\begin{equation}\label{mr:eq:bilog}
 \boxed{\Lambda_{2,(1,1)}(a,b)
 =-2\gamma_1(a)+\sum_{k\ge1}(-h)^k
    \left\{\zeta^{[1]}(k+1,a)+\frac{\zeta(k+1,a)}k\right\}.}
\end{equation}
Only the $k$th Pochhammer of the second shift contributes, with
$\stir{k}{1}=(k-1)!$; its remaining derivative is $E_a^{[1]}(k)$.
This verifies every coefficient and sign in the displayed example.
Although the right side uses the anchor $a$ asymmetrically, its value is
symmetric in $a,b$ after completion, by the defining integral. Expanding
about $b$ therefore gives an exact identity between two Stieltjes--zeta
series whenever their convergence domains overlap.

\section{Balanced log ratios and harmonic-number generating identities}
\label{mr:sec:balanced}
A particularly useful zero-total-order slice avoids a spurious loss of
the constant term. It also gives convergent identities without any
retained Stieltjes constants after one derivative.

\begin{theorem}[Balanced generating trace]\label{mr:thm:balanced}
For $\Rea a,\Rea b>0$ and every $s\in\C$,
\begin{equation}\label{mr:eq:balanced}
 \boxed{\M_2(s,-s;a,b)=1+s(b-a)
     \sum_{n=0}^\infty\frac{(n+b)^{s-1}}{(n+a)^{s+1}}.}
\end{equation}
The series and all its fixed spectral derivatives converge locally
normally throughout that domain. At $a=b$ the right side is one.
\end{theorem}
\begin{proof}
Initially take $a>0$, $h=b-a$, and $|h|<a$. The anchor series gives
$1+\sum_{k\ge1}\binom{s}{k}kh^k\zeta(k+1,a)$, because
$(-1)^k(-s)_k/k!=\binom{s}{k}$ and $E_a(k)=k\zeta(k+1,a)$.
Expanding the zeta series and differentiating the binomial identity
$\sum_{k\ge0}\binom{s}{k}z^k=(1+z)^s$ gives
\eqref{mr:eq:balanced} with absolute interchange in this domain.
For general right-half-plane $a,b$, each summand on its right is
$O(n^{-2})$ locally uniformly, since
$\Log(n+b)-\Log(n+a)=O(n^{-1})$. Spectral derivatives insert powers
of that difference and improve the tail. Thus the right side is jointly
holomorphic in $a,b,s$. The identity theorem on the connected product
of half-planes extends the equality. The branches agree because each
$n+a,n+b$ stays in the right half-plane.
\end{proof}

For real $a,b>1/2$, use the normalized bracket
\begin{equation}\label{mr:eq:bracket}
 \langle f\rangle=\frac\pi2\int_\R\frac{f(t)}{\cosh^2\pi t}\dd t.
\end{equation}
Both resolvents lie in the same right half-plane, so the principal
logarithm of their ratio equals the difference of their principal logs.
Differentiation of \eqref{mr:eq:balanced} gives, for $q\ge1$,
\begin{equation}\label{mr:eq:ratio}
 \boxed{\left\langle
 \Log^q\frac{b-1/2+it}{a-1/2+it}\right\rangle
 =q(b-a)\sum_{n\ge0}
 \frac{\log^{q-1}((n+b)/(n+a))}{(n+a)(n+b)}.}
\end{equation}
In the complex or crossing case the same identity holds for the
completed linear combination of mixed logarithmic moments obtained by
expanding $(L_b-L_a)^q$; no uncompleted branch substitution is intended.

\begin{corollary}[A harmonic coefficient family]\label{mr:cor:harmonic}
For real $a,b>1/2$, $h=b-a$, $|h|<a$, and $q\ge1$,
\begin{equation}\label{mr:eq:harmonic}
 \left\langle\Log^q\frac{b-1/2+it}{a-1/2+it}\right\rangle
 =q!\sum_{k\ge q}(-1)^{k-q}\stir{k}{q}
       \frac{h^k}{(k-1)!}\zeta(k+1,a).
\end{equation}
All parameter differentiations allowed by the preceding theorems may
be made termwise on compact subdisks.
\end{corollary}
\begin{proof}
Differentiate the convergent binomial expansion in the proof of
Theorem~\ref{mr:thm:balanced}, using the signed falling-factorial
coefficient $(-1)^{k-q}\stir{k}{q}$.
\end{proof}
The first three cases are
\begin{align}
 \left\langle\Log\frac{b-1/2+it}{a-1/2+it}\right\rangle
 &=\psi(b)-\psi(a),\label{mr:eq:ratio1}\\
 \left\langle\Log^2\frac{b-1/2+it}{a-1/2+it}\right\rangle
 &=2\sum_{k\ge2}(-h)^kH_{k-1}\zeta(k+1,a),\label{mr:eq:ratio2}\\
 \left\langle\Log^3\frac{b-1/2+it}{a-1/2+it}\right\rangle
 &=-3\sum_{k\ge3}(-h)^k
       \bigl(H_{k-1}^2-H_{k-1}^{(2)}\bigr)\zeta(k+1,a).
 \label{mr:eq:ratio3}
\end{align}
For $q=1$, the series on the right of \eqref{mr:eq:ratio} telescopes
as the standard convergent digamma difference. These identities are
exact analytic generating families, not assertions of reduction of an
arbitrary Euler sum to ordinary zeta products.

\section{Integer orders: Hermite reflection and centered half residues}
\label{mr:sec:integer}
When every order is a positive integer, cuts disappear but residues do
not. Merge repeated branch points and let their multiplicities be $m_l$.
Write
\begin{equation}\label{mr:eq:integerf}
 f(p)=g(p)^R\prod_j(p-p_j)^{-n_j},\qquad n_j\in\N_{>0},
 \qquad N=\sum_jn_j.
\end{equation}

\begin{theorem}[Residue-completed integer reflection]\label{mr:thm:integer}
For $\Rea R>0$ and $\Rea a_j>0$,
\begin{equation}\label{mr:eq:integercentral}
 \M_R(\boldsymbol n;\boldsymbol a)
 =\frac1{2\pi}\Fp\int_\R f(1/2+it)\dd t
   +\sum_{\Rea p_l>1/2}\Res_{p=p_l}f(p)
   +\frac12\sum_{\Rea p_l=1/2}\Res_{p=p_l}f(p).
\end{equation}
The finite part is symmetric Hadamard subtraction at every pole on the
central line, and is an ordinary integral when there are none.
If also $0<\Rea a_j<1$, then
\begin{equation}\label{mr:eq:Hermite}
 \boxed{\M_R(\boldsymbol n;\boldsymbol a)
       -(-1)^N\M_R(\boldsymbol n;\boldsymbol1-\boldsymbol a)
       =\sum_l\Res_{p=p_l}f(p).}
\end{equation}
The right side is the confluent Hermite divided difference of $g^R$
at the multiset containing $p_l$ with multiplicity $m_l$.
\end{theorem}
\begin{proof}
Without central poles, ordinary contour deformation gives the first
formula with the full residues to the right. For a pole on the line,
the boundary distribution computed in Theorem~\ref{mr:thm:contact}
below gives exactly half its residue and the specified finite part.
This can also be obtained by a symmetric indentation and direct
integration of the local Laurent polynomial. Higher order poles
require Hadamard finite parts, not merely principal values.

For the reflected identity replace $p$ by $1-p$ in a right contour
for the second transform. The equality $g(1-p)^R=g(p)^R$ holds with
our logarithm, and the integer denominators supply $(-1)^N$.
The result is a left contour for $f$. Choose the two lines so that all
$p_l$ lie between them; this is possible because $0<\Rea p_l<1$.
Their difference is the sum of residues, proving \eqref{mr:eq:Hermite}.
The usual Cauchy-integral definition of a Hermite divided difference
is precisely this sum. Equivalently, one may expand every local
analytic factor in its Taylor polynomial of degree $m_l-1$, which
gives the standard confluent derivative formula directly.
\end{proof}

For distinct simple shifts the two-factor result is particularly short:
\begin{equation}\label{mr:eq:tworeflect}
 \M_R(1,1;a,b)-\M_R(1,1;1-a,1-b)
 =\frac{g(a)^R-g(b)^R}{b-a}.
\end{equation}
Three distinct simple shifts give
\begin{equation}\label{mr:eq:threereflect}
 \M_R(1,1,1;a,b,c)+\M_R(1,1,1;1-a,1-b,1-c)
 =\sum_{v\in\{a,b,c\}}
   \frac{g(v)^R}{\prod_{u\in\{a,b,c\}\setminus\{v\}}(u-v)}.
\end{equation}
Their limits are Hermite derivatives; apparent small denominators are
not singularities of the identity.

\begin{proposition}[A terminating confluent residue polynomial]\label{mr:prop:Pn}
Let
\begin{equation}\label{mr:eq:Pn}
 P_{0,R}(z)=1,\qquad
 P_{n+1,R}(z)=RzP_{n,R}(z)+(1+z^2)\partial_zP_{n,R}(z).
\end{equation}
At a common shift $a$ with $0<\Rea a<1$,
\begin{equation}\label{mr:eq:commoninteger}
 \M_R(N;a)-(-1)^N\M_R(N;1-a)
 =\frac{\pi^{N-1}g(a)^R}{(N-1)!}P_{N-1,R}(\cot\pi a).
\end{equation}
Every polynomial has rational coefficients in $R,z$. In particular,
$P_{1,R}=Rz$ and $P_{2,R}=R+R(R+1)z^2$.
\end{proposition}
\begin{proof}
At $p=1-a$, $\cot\pi p=-\cot\pi a$. Repeated differentiation of
$g(p)^R$, using $(\cot\pi p)'=-\pi(1+\cot^2\pi p)$, gives
$\partial_p^n g(p)^R|_{p=1-a}=\pi^ng(a)^RP_{n,R}(\cot\pi a)$.
The residue of $g(p)^R/(p-(1-a))^N$ is its $(N-1)$st derivative
divided by $(N-1)!$. Apply \eqref{mr:eq:Hermite}.
\end{proof}
At $a=1/2,N=1$, formula \eqref{mr:eq:integercentral} yields
$\M_R(1;1/2)=\pi^R/2$: the central principal-value integral is zero
by oddness. Dropping the half residue would incorrectly give zero.

\section{All spectral derivatives of the local contact distribution}
\label{mr:sec:contact}
This section fixes the local convention behind all centered integer
collisions. Its distributional formulas are applications of analytic
continuation of one-sided powers; they are proved here to make every
phase and contact coefficient explicit.

Let $t_+^{\lambda}$ and $t_-^{\lambda}$ denote the one-sided powers,
acting initially by integration over $t>0$ and $t<0$. Continue them by
subtracting the Taylor polynomial of a Schwartz test function at zero
on the interval $(0,1)$ in the positive radial coordinate. The cutoff
one fixes a common finite-part scale. For $m\ge1$ define
$A_{m,\ell}^{\pm}$ by the Laurent expansions
\begin{align}
 t_+^{-m-z}
 &=\frac{(-1)^m}{(m-1)!z}\delta^{(m-1)}
   +\sum_{\ell\ge0}\frac{(-z)^\ell}{\ell!}A_{m,\ell}^+,
 \label{mr:eq:halfplus}\\
 t_-^{-m-z}
 &=-\frac1{(m-1)!z}\delta^{(m-1)}
   +\sum_{\ell\ge0}\frac{(-z)^\ell}{\ell!}A_{m,\ell}^-.
 \label{mr:eq:halfminus}
\end{align}
These are the fixed-scale finite parts of
$t_\pm^{-m}\log^\ell|t|$. The series are locally convergent as
meromorphic distribution germs when paired with a Schwartz function.

\begin{theorem}[Entire boundary family and all contact jets]\label{mr:thm:contact}
The boundary value
\begin{equation}\label{mr:eq:T}
 T_s(t)=\lim_{\epsilon\downarrow0}(\epsilon+it)^{-s}
 =e^{-i\pi s/2}t_+^{-s}+e^{i\pi s/2}t_-^{-s}
\end{equation}
is an entire tempered-distribution family in $s$. Near $s=m+z$,
\begin{align}
 T_{m+z}={}&e^{-i\pi(m+z)/2}
       \sum_{\ell\ge0}\frac{(-z)^\ell}{\ell!}A_{m,\ell}^+
 +e^{i\pi(m+z)/2}
       \sum_{\ell\ge0}\frac{(-z)^\ell}{\ell!}A_{m,\ell}^-\nonumber\\
 &+\frac{2i^{m-1}}{(m-1)!}\frac{\sin(\pi z/2)}z\,
       \delta^{(m-1)}.
 \label{mr:eq:contactgerm}
\end{align}
In this fixed decomposition, the coefficient of $\delta^{(m-1)}$ in
the $q$th order derivative at $s=m$ is
\begin{equation}\label{mr:eq:contactcoef}
 c_{m,q}=\begin{cases}
 \displaystyle\frac{\pi i^{m-1}}{(m-1)!}
      \frac{(-1)^{q/2}(\pi/2)^q}{q+1},&q\text{ even},\\[5pt]
 0,&q\text{ odd}.
 \end{cases}
\end{equation}
The nonlocal part is the finite explicit sum
\begin{align}
 &e^{-i\pi m/2}\sum_{\ell=0}^q\binom q\ell
       (-i\pi/2)^{q-\ell}(-1)^\ell A_{m,\ell}^+\nonumber\\
 &\qquad+e^{i\pi m/2}\sum_{\ell=0}^q\binom q\ell
       (i\pi/2)^{q-\ell}(-1)^\ell A_{m,\ell}^-.
 \label{mr:eq:contactnonlocal}
\end{align}
\end{theorem}
\begin{proof}
For $\Rea s<1$, the right side of \eqref{mr:eq:T} is locally
integrable and is the ordinary limit of $(\epsilon+it)^{-s}$.
Taylor subtraction continues its action on Schwartz functions
meromorphically to the plane. The only possible poles are positive
integers, with residues obtained by integrating the $(m-1)$st Taylor
term. They are exactly \eqref{mr:eq:halfplus}--\eqref{mr:eq:halfminus}.
After multiplication by the phases their polar contribution is
\[
 \frac{i^m}{(m-1)!z}
     (e^{-i\pi z/2}-e^{i\pi z/2})\delta^{(m-1)}
 =\frac{2i^{m-1}}{(m-1)!}\frac{\sin(\pi z/2)}z\delta^{(m-1)}.
\]
It is entire at $z=0$, so all poles cancel. The exponential-boundary
regularization has the same continuation: subtract the test-function
Taylor terms first and integrate those elementary terms explicitly.
This establishes its distributional limit and the entire identity.
At integer orders there is also a direct check: the real and imaginary
parts of $(\epsilon+it)^{-1}$ are
$\epsilon/(\epsilon^2+t^2)$ and $-t/(\epsilon^2+t^2)$, whose limits
are $\pi\delta$ and $-\operatorname{PV}(1/t)$. Repeated
$t$ differentiation gives all positive integer boundary values.
For noninteger orders, integration by parts reduces the limit to a
locally integrable order; Taylor subtraction makes the reduction
locally uniform through the removable integer parameters.
Taking Taylor coefficients of the sine quotient proves
\eqref{mr:eq:contactcoef}; Leibniz's rule on the two regular series gives
\eqref{mr:eq:contactnonlocal}.
\end{proof}
At $q=0$, this reduces to
\begin{equation}\label{mr:eq:plemelj}
 T_m=i^{-m}\Fp t^{-m}
       +\frac{\pi i^{m-1}}{(m-1)!}\delta^{(m-1)}.
\end{equation}
To recover the half residue in \eqref{mr:eq:integercentral}, multiply
by a local analytic factor $H(1/2+it)$ and integrate with $1/(2\pi)$.
The delta derivative contributes
$H^{(m-1)}(p_0)/(2(m-1)!)$, since the two powers of $i^{m-1}$ and
the distributional sign cancel. This is precisely half the residue.
Translation gives the same rule at $p_0=1/2+iy_0$.

\begin{remark}[Canonical total, conventional splitting]
Changing the finite-part scale redistributes local delta derivatives
between the $A_{m,\ell}^{\pm}$ and the displayed contact coefficients.
The full distribution $T_s$ and every derivative are canonical.
Formula \eqref{mr:eq:contactcoef} is stated only with the common scale
fixed above. It must not be combined with finite parts taken at a
different scale. If several centered factors have the same point,
merge their orders first; independent spectral derivatives at that
collision become derivatives in their total order. Before collision
all orders remain independent in \eqref{mr:eq:master}.
\end{remark}

\section{Centered Fourier--Mellin reflection and the entire alternating family}
\label{mr:sec:centered}
The second research continuation requires the full spectral family,
not just its integer values. Define
\begin{equation}\label{mr:eq:Z}
 Z_R(v)=\frac1{\Gamma(v)}\int_0^\infty
 t^{v-1}(2\cosh(t/2))^{-R}\dd t,\qquad\Rea v>0,\quad\Rea R>0.
\end{equation}
Its original absolutely convergent alternating series is
\begin{equation}\label{mr:eq:Zseries}
 Z_R(v)=\sum_{n\ge0}\frac{(-1)^n(R)_n}{n!(n+R/2)^v},
 \qquad\Rea v>\Rea R.
\end{equation}
The generalized series is not an ordinary Lerch function at arbitrary
$R$. It is useful to retain its explicit definition.

\begin{lemma}[Entire continuation and centered reflection]\label{mr:lem:reflection}
The function $Z_R(v)$ is entire in $v$ and jointly holomorphic for
$\Rea R>0$. If $f_R(t)=(2\cosh(t/2))^{-R}$, then
\begin{equation}\label{mr:eq:negativeZ}
 Z_R(-n)=(-1)^n f_R^{(n)}(0),\qquad
 Z_R(-2m-1)=0.
\end{equation}
Writing $F_R(w)=\F_{R,1/2}(w;0)$, one has
\begin{equation}\label{mr:eq:reflection}
 F_R(w)=2(2\pi)^{R+w-2}\Gamma(1-w)
              \cos(\pi w/2)Z_R(1-w),
\end{equation}
with all apparent singularities removable.
\end{lemma}
\begin{proof}
Subtract a Taylor polynomial of $f_R$ at zero in the integral over
$(0,1)$. Its Mellin transform has simple poles at $v=-n$ with residue
$f_R^{(n)}(0)/n!$. The reciprocal Gamma factor cancels those poles and
has derivative $(-1)^n n!$ there, proving \eqref{mr:eq:negativeZ}.
Evenness of $f_R$ gives the odd zeros. Locally uniform subtractions and
exponential decay at infinity prove joint holomorphy. The binomial
series gives \eqref{mr:eq:Zseries} in its absolute region.

For $0<\Rea w<1$, the fractional formula gives
$F_R(w)=\Gamma(w)^{-1}\int_0^\infty x^{w-1}b_R(x)\dd x$.
By \eqref{mr:eq:beta}, $b_R(x)=(2\pi)^{R-2}\widehat f_R(x/(2\pi))$
with $\widehat f(v)=\int_\R e^{ivt}f(t)\dd t$.
The Mellin transform of the cosine Fourier transform is
\[
 \int_0^\infty v^{w-1}\widehat f_R(v)\dd v
 =2\Gamma(w)\cos(\pi w/2)\int_0^\infty t^{-w}f_R(t)\dd t.
\]
For completeness, insert $e^{-\epsilon v}$ first. Fubini is then
absolute; the inner integral is
$\Gamma(w)((\epsilon-it)^{-w}+(\epsilon+it)^{-w})/2$.
It is uniformly dominated by $Ct^{-\Rea w}$ for $t>0$, integrable
against $f_R$ in this strip. Removing the regulator and rescaling
proves \eqref{mr:eq:reflection}. Entirety of $F_R$, supplied by its
right contour, proves the removability and the global identity.
\end{proof}
These inputs are established in the continuous-depth report \cite{CDL};
the preceding proof records the normalization used by the new jet formulas.
In particular,
\begin{equation}\label{mr:eq:classicalZ}
 Z_1(v)=2^v\beta(v),\qquad Z_2(v)=\eta(v-1),
\end{equation}
where $\beta$ is the Dirichlet beta function and
$\eta(v)=(1-2^{1-v})\zeta(v)$ is entire.

\section{All resonant spectral jets: finite Bell formulas}
\label{mr:sec:resonance}
The following theorem answers the spectral-jet question without
illegitimately differentiating an identity known only on a discrete set.
Set $L=\log(2\pi)$. For $n\ge0$ and $k\ge1$ define two cumulant lists:
\begin{align}
 d^{\mathrm o}_{n,1}=d^{\mathrm e}_{n,1}
 &=L+\gamma-H_n,\label{mr:eq:d1}\\
 d^{\mathrm o}_{n,k}
 &=(k-1)!\left\{(1-2^{1-k}\mathbf1_{2\mid k})\zeta(k)
                      +(-1)^k H_n^{(k)}\right\},\quad k\ge2,
 \label{mr:eq:do}\\
 d^{\mathrm e}_{n,k}
 &=(k-1)!\left\{(1-2(1-2^{-k})\mathbf1_{2\mid k})\zeta(k)
                      +(-1)^k H_n^{(k)}\right\},\quad k\ge2.
 \label{mr:eq:de}
\end{align}
Let $B_{n,q}^{\mathrm o}=\B_q(d^{\mathrm o}_{n,1},\ldots,d^{\mathrm o}_{n,q})$
and define $B_{n,q}^{\mathrm e}$ similarly, with both zero-index values one.

\begin{theorem}[Every odd and even order jet]\label{mr:thm:resonance}
For $\Rea R>0$, $q\ge0$, and $m\ge0$,
\begin{align}
 F_R^{[q]}(2m+1)
 ={}&\frac{(-1)^m\pi(2\pi)^{R+2m-1}}{(2m)!}
 \sum_{j=0}^q\binom qj B_{2m,q-j}^{\mathrm o}(-1)^j
                           Z_R^{[j]}(-2m).
 \label{mr:eq:oddjet}
\end{align}
For $m\ge1$,
\begin{align}
 F_R^{[q]}(2m)
 ={}&\frac{2(-1)^{m+1}(2\pi)^{R+2m-2}}{(2m-1)!}
 \sum_{j=0}^q\binom qj B_{2m-1,q-j}^{\mathrm e}
           \frac{(-1)^j}{j+1}Z_R^{[j+1]}(1-2m).
 \label{mr:eq:evenjet}
\end{align}
Both sums are finite. All mixed $R$ derivatives are also valid.
\end{theorem}
\begin{proof}
At the odd point $w=2m+1+z$, Gamma recurrence in
\eqref{mr:eq:reflection} gives the analytic germ
\begin{equation}\label{mr:eq:oddgerm}
 F_R(2m+1+z)=C_{R,m}^{\mathrm o}A_{2m}^{\mathrm o}(z)Z_R(-2m-z),
\end{equation}
where
\begin{align}
 C_{R,m}^{\mathrm o}&=\frac{(-1)^m\pi(2\pi)^{R+2m-1}}{(2m)!},\nonumber\\
 A_n^{\mathrm o}(z)&=(2\pi)^z\Gamma(1-z)
     \frac{\sin(\pi z/2)}{\pi z/2}
     \prod_{k=1}^n(1+z/k)^{-1}.\label{mr:eq:Ao}
\end{align}
At the even point, put $n=2m-1$ and use $Z_R(-n)=0$. The analytic germ is
\begin{equation}\label{mr:eq:evengerm}
 F_R(2m+z)=C_{R,m}^{\mathrm e}A_n^{\mathrm e}(z)
                         \left(-\frac{Z_R(-n-z)}z\right),
\end{equation}
where
\begin{align}
 C_{R,m}^{\mathrm e}&=\frac{2(-1)^{m+1}(2\pi)^{R+2m-2}}{(2m-1)!},\nonumber\\
 A_n^{\mathrm e}(z)&=(2\pi)^z\Gamma(1-z)\cos(\pi z/2)
                    \prod_{k=1}^n(1+z/k)^{-1}.\label{mr:eq:Ae}
\end{align}
All factors have value one at zero except the explicitly retained $Z_R$
terms. The convergent log-Gamma expansion \cite{GammaSeries}, the
sine and cosine products, and the finite logarithmic product give
\[
 \log A_n^{\mathrm o}(z)=\sum_{k\ge1}d^{\mathrm o}_{n,k}\frac{z^k}{k!},
 \qquad
 \log A_n^{\mathrm e}(z)=\sum_{k\ge1}d^{\mathrm e}_{n,k}\frac{z^k}{k!}.
\]
For example, the logarithm of the sinc factor contributes
$-2^{1-k}\zeta(k)/k$ at each even degree $k$; that of the cosine
contributes $-2(1-2^{-k})\zeta(k)/k$. The finite product contributes
$(-1)^kH_n^{(k)}/k$. These are exactly
\eqref{mr:eq:d1}--\eqref{mr:eq:de}.
Exponentiation gives the Bell coefficients. Finally,
\[
 -\frac{Z_R(-n-z)}z
 =\sum_{j\ge0}\frac{(-1)^jZ_R^{[j+1]}(-n)}{(j+1)!}z^j.
\]
The product rule proves both asserted sums, including the factor
$1/(j+1)$ in the even case. Joint holomorphy justifies $R$ derivatives.
\end{proof}

For clarity, the first odd derivative is
\begin{equation}\label{mr:eq:odd1}
 F_R^{[1]}(2m+1)=C_{R,m}^{\mathrm o}
 \left\{(L+\gamma-H_{2m})Z_R(-2m)-Z_R^{[1]}(-2m)\right\}.
\end{equation}
The first even derivative is
\begin{equation}\label{mr:eq:even1}
 F_R^{[1]}(2m)=C_{R,m}^{\mathrm e}
 \left\{(L+\gamma-H_{2m-1})Z_R^{[1]}(1-2m)
                       -\tfrac12Z_R^{[2]}(1-2m)\right\}.
\end{equation}
No division by a possibly vanishing $Z_R$ value occurs.
At $q=0$ the odd formula recovers the prior polynomial evaluation
\begin{equation}\label{mr:eq:oddvalue}
 F_R(2m+1)=\frac{\pi^{R+2m}}2U_m(R),\qquad
 U_m(R)=[z^{2m}]\sec(z)^R.
\end{equation}
The jets $Z_R^{[j]}(-2m)$ for $j>0$ are genuinely retained spectral
data. Knowing \eqref{mr:eq:oddvalue} for all integer $m$ is not, by
itself, a license to differentiate it in $m$.

\section{Gamma-weighted polylogarithms and explicit Stieltjes order jets}
\label{mr:sec:polylog}
For real $x$, use principal polylogarithms and put
\begin{equation}\label{mr:eq:polyodd}
 \mathcal P_s(x)=\frac{\Li_s(ie^{-x/2})-\Li_s(-ie^{-x/2})}{i},
 \qquad I_R(s)=\int_\R G_R(x)\mathcal P_s(x)\dd x.
\end{equation}
Their arguments do not cross the positive-real polylogarithm cut.

\begin{theorem}[Entire-order Gamma--polylogarithm reflection]\label{mr:thm:polylog}
For $\Rea R>0$ and every complex $s$, the integral defining $I_R(s)$
is ordinary and absolutely convergent, and
\begin{align}
 I_R(s)&=\Gamma(R)2^{1-s}(2\pi)^{2-R}F_{R+1}(s),
 \label{mr:eq:ItoF}\\
 &=8\Gamma(R)\pi^{s+1}\Gamma(1-s)
           \cos(\pi s/2)Z_{R+1}(1-s).
 \label{mr:eq:Ireflection}
\end{align}
All apparent singularities are removable. Every fixed mixed spectral
and depth derivative can be taken under the ordinary integral.
\end{theorem}
\begin{proof}
Separating the odd terms of the polylogarithm series proves
\[
 \F_{1,1/2}(s;-x)=2^{s-1}e^{-x/2}\mathcal P_s(x)
\]
first in its series region and then along the real axis by analytic
continuation. The Lerch-to-polylogarithm convention is fixed by the
standard relation $\Li_s(z)=z\Phi(z,s,1)$ \cite{Lerch,Polylog}.
Convolve this function with $K_R$ at zero. The common-shift convolution
in Theorem~\ref{mr:thm:master} gives $F_{R+1}(s)$, while the
Gamma factor in \eqref{mr:eq:kernel} gives \eqref{mr:eq:ItoF}.
Its weighted Schwartz estimates already prove absolute convergence
and all mixed differentiations. Substituting \eqref{mr:eq:reflection}
and collecting powers of $2$ and $\pi$ yields
\eqref{mr:eq:Ireflection}.
\end{proof}

\subsection{Odd and even spectral derivatives}
Theorem~\ref{mr:thm:resonance} applies without another singular limiting
calculation. In its cumulants replace $L=\log(2\pi)$ by $\log\pi$,
replace $Z_R$ by $Z_{R+1}$, and replace the two prefactors by
\begin{equation}\label{mr:eq:Ipre}
 \widetilde C_{R,m}^{\mathrm o}
 =\frac{4(-1)^m\Gamma(R)\pi^{2m+3}}{(2m)!},\qquad
 \widetilde C_{R,m}^{\mathrm e}
 =\frac{8(-1)^{m+1}\Gamma(R)\pi^{2m+1}}{(2m-1)!}.
\end{equation}
The resulting finite sums evaluate $I_R^{[q]}(2m+1)$ and
$I_R^{[q]}(2m)$ respectively. The shift in $L$ is necessary because
of the $2^{-s}$ factor in \eqref{mr:eq:ItoF}.

At zero derivative order one recovers the prior odd family
\begin{equation}\label{mr:eq:Iodd}
 I_R(2m+1)=\frac{\Gamma(R)\pi^{2m+3}}{2^{R+2m-1}}U_m(R+1),
\end{equation}
and obtains the even counterpart
\begin{equation}\label{mr:eq:Ieven}
 I_R(2m)=\frac{8(-1)^{m+1}\Gamma(R)\pi^{2m+1}}{(2m-1)!}
                  Z_{R+1}^{[1]}(1-2m),\qquad m\ge1.
\end{equation}
The latter is an exact spectral-derivative coordinate, not an assertion
that it is a polynomial in $R$ like the odd values.

\subsection{Two completely classical depth reductions}
The cases $R=1,2$ provide useful fully explicit families:
\begin{align}
 I_1(s)&=4\pi s(2-2^{-s})\zeta(s+1),\label{mr:eq:I1}\\
 I_2(s)&=4s(s+1)\beta(s+2)+\pi^2\beta(s).
 \label{mr:eq:I2}
\end{align}
These formulas are entire identities, with removable values understood.
For \eqref{mr:eq:I1}, use $F_2(s)=s\zeta(s+1,1/2)$ in
\eqref{mr:eq:ItoF}, and
$\zeta(v,1/2)=(2^v-1)\zeta(v)$.
For \eqref{mr:eq:I2}, Gamma recurrence gives
$K_3(x)=(x^2+\pi^2)K_1(x)/2$. The fractional formula at the center
therefore yields
\[
 F_3(s)=\tfrac12\{s(s+1)F_1(s+2)+\pi^2F_1(s)\},
 \qquad F_1(s)=2^s\beta(s).
\]
Substitution proves \eqref{mr:eq:I2} first at positive real orders
and then everywhere. These are consequences of classical recurrences;
we do not assign separate priority to their low-order values.

Nevertheless their spectral derivatives give concrete ordinary integrals
of polylogarithm order derivatives. In particular,
\begin{align}
 I_1(0)&=4\pi,\nonumber\\
 I_1^{[1]}(0)&=4\pi(\gamma+\log2),\label{mr:eq:I1jets}\\
 I_1^{[2]}(0)&=4\pi\{-2\gamma_1+2\gamma\log2-\log^22\},\nonumber\\
 I_1^{[3]}(0)&=4\pi\{3\gamma_2-6\gamma_1\log2
                          -3\gamma\log^22+\log^32\}.\nonumber
\end{align}
Here $I_1^{[q]}(0)$ means the ordinary integral in
\eqref{mr:eq:polyodd} with $\partial_s^q\mathcal P_s|_{s=0}$;
$G_1(x)=\pi/\cosh(x/2)$. Thus these identities link ordinary zeta
Laurent coefficients to genuine polylogarithm order derivatives,
not to formal derivatives of a discrete family.
For every $q$, a finite closed coefficient rule is
\begin{equation}\label{mr:eq:I1all}
 I_1^{[q]}(0)=4\pi\sum_{j=0}^q\binom qj a_j e_{q-j},\quad
 \begin{cases}
 a_0=e_0=1,\\
 a_j=(-1)^{j+1}\log^j2 &(j\ge1),\\
 e_k=k(-1)^{k-1}\gamma_{k-1} &(k\ge1).
 \end{cases}
\end{equation}
It follows by expanding $(2-2^{-s})E_1(s)$, and agrees with all the
displayed cases.

\subsection{Polygamma-weighted versions}
Define
\begin{equation}\label{mr:eq:depthweights}
 D_k(R,x)=2^{-k}\left\{
 \psi^{(k-1)}\left(\frac R2+\frac{ix}{2\pi}\right)
 +\psi^{(k-1)}\left(\frac R2-\frac{ix}{2\pi}\right)\right\}.
\end{equation}
Logarithmic differentiation of the Gamma product gives
\begin{equation}\label{mr:eq:depthjet}
 \partial_R^d\partial_s^q I_R(s)=
 \int_\R G_R(x)\B_d(D_1(R,x),\ldots,D_d(R,x))
                   \partial_s^q\mathcal P_s(x)\dd x.
\end{equation}
The right side is ordinary for every fixed $d,q$.
Combining it with \eqref{mr:eq:Ireflection}, or with the resonance
formulas before taking the derivative, gives all mixed polygamma and
polylogarithm-order identities. The $D_k$ are genuine Gamma-argument
polygammas; the resulting $Z$ depth derivatives remain distinct
parameters and are not silently discarded.

\section{All anchored primitive orders and an explicit log-Gamma integral}
\label{mr:sec:primitives}
The independent argument ladder follows directly from the defining contour:
\begin{equation}\label{mr:eq:shiftladder}
 \partial_{a_j}\M_R(\boldsymbol s;\boldsymbol a)
 =-s_j\M_R(\boldsymbol s+\boldsymbol e_j;\boldsymbol a).
\end{equation}
The scalar primitive ladder and its removable resonances are already
part of the incoming Lerch calculus \cite{CDL,Tilt}. The following
form makes the independent-shift completion and its normalization explicit.

Let $\mathcal I^k_{A,a_j}$ denote the $k$-fold primitive in $a_j$
with its first $k$ initial derivatives zero at the right-half-plane
anchor $A$. For a holomorphic function, integration may be taken on
the straight segment; other paths inside the half-plane give the same
anchored primitive.

\begin{proposition}[Independent anchored primitives]\label{mr:prop:primitive}
Put
$H(\boldsymbol a)=\M_R(\boldsymbol s-k\boldsymbol e_j;\boldsymbol a)$.
Then
\begin{equation}\label{mr:eq:primitive}
 \mathcal I^k_{A,a_j}\M_R(\boldsymbol s;\boldsymbol a)
 =\frac{H(\boldsymbol a)-
       \displaystyle\sum_{\ell=0}^{k-1}\frac{(a_j-A)^\ell}{\ell!}
         \left.\partial_{a_j}^{\ell}H\right|_{a_j=A}}
       {(1-s_j)_k}.
\end{equation}
The apparent singularities at $s_j=1,\ldots,k$ are removable.
All fixed mixed order, depth, and remaining-shift derivatives are valid.
\end{proposition}
\begin{proof}
Repeated use of \eqref{mr:eq:shiftladder} gives
$\partial_{a_j}^kH=(-1)^k(s_j-k)_k\M_R=(1-s_j)_k\M_R$.
Taylor's formula with integral remainder proves \eqref{mr:eq:primitive}
where the denominator is nonzero. Its right side has the same anchored
integral representation after cancellation, which proves holomorphic
removability at every resonance. Each zero of $(1-s_j)_k$ is simple;
there the quotient may also be computed as the spectral derivative
of the full numerator divided by the derivative of the denominator.
The integral remainder and compact-set holomorphy justify all the
stated mixed derivatives.
\end{proof}

The scalar Stieltjes primitive is worth recording separately:
\begin{equation}\label{mr:eq:stieltjesprimitive}
 \int_A^a\gamma_m(u)\dd u
 =\frac{(-1)^{m+1}}{m+1}
   \{\zeta^{[m+1]}(0,a)-\zeta^{[m+1]}(0,A)\}.
\end{equation}
Indeed $\partial_a\zeta(v,a)=-v\zeta(v+1,a)$ \cite{Zeta};
comparison with \eqref{mr:eq:stieltjes} proves the formula including
its sign. At $m=0$ it is
$\int_A^a-\psi(u)\dd u=-\log\Gamma(a)+\log\Gamma(A)$.

An ordinary completed Gamma primitive follows directly from the raw
log moment. For real $a>0$, let
\begin{equation}\label{mr:eq:VB}
 V(z)=z(\Log z-1),\quad V(0)=0,\qquad
 B(a)=\begin{cases}\log\sin(\pi a),&0<a<1/2,\\0,&a\ge1/2.\end{cases}
\end{equation}
For all real $a,A>0$,
\begin{equation}\label{mr:eq:gammaprimitive}
 \boxed{\begin{aligned}
 &\frac\pi2\int_\R\frac{
        V(a-1/2+it)-V(A-1/2+it)}{\cosh^2\pi t}\dd t
       -B(a)+B(A)\\
 &\hspace{35mm}=\log\Gamma(a)-\log\Gamma(A).
 \end{aligned}}
\end{equation}
To prove it, integrate \eqref{mr:eq:digamma} between $A$ and $a$.
Fubini is justified by the integrable logarithmic singularity and the
exponential kernel. Since $V'=\Log$ away from the cut, continuity at
the center gives the same integral across it. The correction is exactly
$B'(a)=\pi\cot\pi a$ below the center, with no jump in $B$ at $1/2$.
Both sides vanish at $a=A$, fixing the integration constant.

\section{Source audit, corrections, and evidentiary status}
\label{mr:sec:audit}
\subsection{Inspected scope}
The repository read was anchored to \pin, whose recorded commit time is
\src{2026-10-11T03:35:18Z} (10 October in the stated Pacific time zone).
The canonical README, chapter listing, incoming directory, the canonical
$S_4$ proof source, and selected research-question search results were
inspected. The relevant sections of the continuous-depth archive's
\src{article.tex} were read: principal definitions, kernel and entire
spectral family, centered reflection, harmonic master, and the complete
further-research section. The READMEs of the tilted-Lerch, mixed-resolvent,
and algebraic-Gauss incoming packages were read to delimit overlap.
Those README reads are \emph{not} full analytic audits of their articles.
The other incoming ZIPs were inventoried, not fully unpacked or certified.
The machine-readable source record names these exact scopes and observed
Git blob identifiers.

The accessible parsed text of Johansson--Blagouchine's original PDF
confirmed the $\Rea a>1/2$ hypothesis of its equation (2).
The web renderer failed to supply the requested screenshot; no visual
inspection of that source PDF is claimed. The present article's PDF
was separately compiled and visually inspected as recorded in the
package's build and visual-review receipts.

\subsection{What is corrected, and what is not}
No false theorem was identified in the inspected source portions.
Instead, three concrete \emph{invalid extensions} are excluded by
explicit counterexamples: removing the crossing correction in
\eqref{mr:eq:digamma}; adding only pairwise logarithmic corrections
and losing the overlap term in \eqref{mr:eq:tripleP}; and dropping
half residues at centered integer poles. A fourth structural safeguard
is to combine a Gamma pole and a trivial $Z_R$ zero before taking
spectral derivatives, as in \eqref{mr:eq:evengerm}.

Additional integration safeguards are precise rather than rhetorical.
One must not replace the complex Gamma product by a squared modulus;
one must use the upper-minus-lower difference instead of an imaginary
part for complex shifts; one must retain the constant $1$ on the
zero-total-order slice \eqref{mr:eq:balanced}; and one must specify the
finite-part scale when isolating the contact coefficients.
A convergence hypothesis on a cut integral does not disappear merely
because its sine prefactor vanishes at an integer order.

The canonical $S_6$ and revised $S_8$ candidates remain unproved here.
The existing $S_4$ certificate is not recounted as a new result.
Neither the retained $Z_R$ jets nor the Hurwitz-zeta derivative series
are claimed to be arithmetically independent or irreducible.

\subsection{Executed checks and what they establish}
The exact script executed \textbf{490 symbolic assertions}, including
three negative controls. They compare Pochhammer derivatives with
Stirling coefficients, mixed jump products with their independent
multinomial expansion, residue recurrences with a trigonometric
translation series, both lists of resonance cumulants with analytic
Taylor germs, contact coefficients with the sine quotient, and
primitive denominators at every tested resonance.

The independent numerical suite passed \textbf{77 positive comparisons
and three negative controls at each of 45 and 60 working decimal digits}.
The largest normalized positive discrepancies were approximately
$5.185\times10^{-31}$ and $7.510\times10^{-41}$, respectively.
The normalization is $|u-v|/(1+|u|+|v|)$.
Tests include unequal and coincident real cuts, complex shifts at
separate and common heights, genuinely complex orders and depth,
ordinary Stieltjes moments, cubic overlap corrections, Hermite residues,
centered half residues, harmonic and mixed-zeta expansions, independent
right-contour spectral jets, Gamma primitives crossing the center, and
Gamma--polylogarithm quadrature.

These are ordinary \texttt{mpmath} floating-point diagnostics, not
outward-rounded interval enclosures. The harmonic checks truncate
at $k=79$; their analytically proved convergence is not a numerical
certificate of the displayed infinite sums. Polylogarithm quadrature
uses finite cutoffs $310$ and $400$ in the two runs, to avoid pathological
huge-argument evaluations at transformed infinite quadrature nodes.
The complete exact identities are proved above, not established by
small numerical residuals. Near spectral zero, the tiny-step Hurwitz
implementation can suffer internal cancellation; the independent
$Z_1$ derivative check uses Cauchy differentiation on a small complex
circle instead. The executed code records these choices.

\section{Further research questions and proposed next targets}
\label{mr:sec:future}
The following are open directions, not additional claims of proof.
They focus on exact identities and structural reductions.

\begin{enumerate}[label=Q\arabic*.]
\item \textbf{Arithmetic closure of the retained spectral jets.}
For rational depth $R$, determine which combinations of
$Z_R^{[j]}(-n)$ reduce to ordinary Hurwitz-zeta derivatives, Gamma
values, or cyclotomic multiple polylogarithms. Start with the even
Gamma--polylogarithm values in \eqref{mr:eq:Ieven}; a formula that
merely renames the same jet is not the desired arithmetic reduction.

\item \textbf{Depth jets at an integer total depth.}
Combine the mixed $R$ derivatives of the present reflection with a
canonical logarithmically extended Lerch basis. The multiplier
$g(p)^N\Log^q g(p)$ is outside the old meromorphic finite-depth class;
find an exact basis with explicit connection and distribution laws.

\item \textbf{Independent Gamma weights and several polylogarithms.}
The present convolution has one polylogarithm odd part against one
Gamma product. Several independent polylogarithms in the same variable
are not handled by simply adding their orders. Determine the multiple
Hurwitz--Lerch objects and balanced combinations that admit a genuine
Gamma-quotient generating function.

\item \textbf{A canonical global mixed-log coordinate system.}
The anchor compiler converges in its natural polydisk; overlapping
anchors give exact transition identities such as the symmetry of
\eqref{mr:eq:bilog}. Find a terminating normal form for those
transitions without assuming linear independence of their special values.

\item \textbf{Rational shift summations of balanced log ratios.}
Classify residue-class averages of \eqref{mr:eq:ratio} that telescope
to finite Gamma or zeta expressions. The shifted harmonic coefficients
in \eqref{mr:eq:harmonic} provide an explicit generating target at
every degree, rather than a numerical integer-relation guess.

\item \textbf{Distributional depth continuation.}
Continue the \emph{depth} $R$ through zero and negative real parts,
including the Gamma poles that cross its argument contour. The contact
analysis in Section~\ref{mr:sec:contact} treats spectral powers, not
this separate depth continuation; it does not authorize discarding
those additional poles.

\item \textbf{A functional certificate for the Gaussian candidates.}
Use the multishift identities to search for convergent mixed-color
relations relevant to $S_6$ and revised $S_8$. A successful continuation
must exhibit the exact relation or a replayable rational certificate;
the present reflection alone does not prove either candidate.

\item \textbf{Formal verification and certified replay.}
Formalize the finite Stirling, Bell, and contact algebra first, then isolate
weighted Fourier inversion, parameter holomorphy, contour deformation,
and distributional Taylor subtraction as analytic obligations.
An independent ball-arithmetic replay would give certified diagnostic
enclosures, but should remain distinct from the arbitrary-index proofs.

\item \textbf{Literature reconciliation at the level of full families.}
Compare the exact multishift jump theorem and its spectral germs with
classical Mellin--Barnes, generalized-Lerch, and multiple-zeta literature.
Assign historical priority only after comparing whole generating
functions and their branch normalizations, not isolated low-weight values.
\end{enumerate}

\section*{Conclusion}
\addcontentsline{toc}{section}{Conclusion}
The independent-shift reflection is governed by the jump of a
\emph{product}, not the sum of isolated jumps. With that correction,
complex shifts, coincident cuts, ordinary mixed logarithmic jets,
and integer half-residue contacts fit one holomorphic transform.
The Stieltjes compiler and harmonic ratio identities give concrete
exact consequences at every degree. Separately, the analytic germs
at both resonance lattices determine all spectral derivatives with
finite Bell coefficients, retaining the spectral data that discrete
value formulas cannot supply. These results continue two specified
repository questions while leaving arithmetic reductions, depth-zero
continuation, and the Gaussian candidates explicitly open.

\appendix
\section{Replay and integration map}
\label{mr:app:replay}
The self-contained source \src{article.tex} builds without a ProveIt
checkout. From the package directory:
\begin{verbatim}
python code/verify_exact.py
python code/verify_numeric.py --dps 45
python code/verify_numeric.py --dps 60
python code/build.py
\end{verbatim}
The build script runs three serial pdfLaTeX passes and checks for
unresolved references and overfull boxes. The JSON receipts retain
working precision, values on both sides, finite cutoffs, and normalized
residuals. The symbolic suite needs SymPy; the numerical suite needs
mpmath; PDF inspection needs PyMuPDF and Pillow. Versions are recorded
in \src{requirements.txt} and the receipts.

The proposed destination is
\src{Analysis/Polylogarithms/docs/reports/multishift-lerch-reflection/}.
The \src{integration/} directory supplies an additive insertion,
claim ledger, and source record. Integration should cite the older
Gamma transform, continuous-depth definitions, and anchor expansion,
then add the completed-cut theorem and resonance-jet theorems after
those prerequisites. No original archive needs to be overwritten.
The canonical volumes were not rebuilt and no repository mutation was
performed in preparing this package.

\begin{thebibliography}{99}
\addcontentsline{toc}{section}{References}
\bibitem{Repo}
V. Reshetnikov and ProveIt contributors,
\emph{ProveIt}, canonical polylogarithm manuscript
and incoming research queue, repository snapshot \pin.
\url{https://github.com/VladimirReshetnikov/ProveIt}.

\bibitem{CDL}
\emph{Continuous-Depth Lerch Calculus: Gamma kernels, all-order harmonic
identities, Stieltjes convolutions, and anchored primitives}, research
continuation prepared for V. Reshetnikov, 10 October 2026.
Incoming archive \src{ProveIt_Continuous_Depth_Lerch_2026-10-10.zip},
member \src{ProveIt_Continuous_Depth_Lerch_2026-10-10/article.tex}.
Read portions and archive Git blob are recorded in the source audit;
especially the centered-reflection section and further-research Q4, Q8.

\bibitem{Tilt}
\emph{Tilted Lerch Convolutions and Binary Harmonic Cancellation},
research continuation prepared for V. Reshetnikov, 10 October 2026.
Incoming archive \src{ProveIt_Tilted_Lerch_Binary_Cancellation_2026-10-10.zip}.
Its README describes the independent-shift single-anchor expansion,
all order jets, and anchored primitives; those inputs are reproved
in the narrower notation needed here rather than assumed from an unread proof.

\bibitem{CS}
D. Chakrabarti and G. K. Srinivasan,
\emph{On a remarkable formula of Ramanujan}, 2012,
arXiv:1206.5122. \url{https://arxiv.org/abs/1206.5122}.

\bibitem{JB}
F. Johansson and I. V. Blagouchine,
\emph{Computing Stieltjes constants using complex integration},
arXiv:1804.01679v3, 13 August 2018, especially equation (2).
\url{https://arxiv.org/abs/1804.01679}.

\bibitem{CH}
T.-L. Cheng and C.-Y. Hu,
\emph{On the Hurwitz Zeta Function and Its Applications to Hyperbolic
Probability Distributions}, 2022, arXiv:2208.13733.
\url{https://arxiv.org/abs/2208.13733}.
This provides additional classical-context comparison for hyperbolic
Hurwitz representations; no uninspected theorem from it is needed here.

\bibitem{Beta}
NIST Digital Library of Mathematical Functions,
Section 5.12, \emph{Beta Function}.
\url{https://dlmf.nist.gov/5.12}.

\bibitem{GammaSeries}
NIST Digital Library of Mathematical Functions,
Section 5.7, \emph{Series Expansions} of the Gamma function.
\url{https://dlmf.nist.gov/5.7}.

\bibitem{Polygamma}
NIST Digital Library of Mathematical Functions,
Section 5.15, \emph{Polygamma Functions}.
\url{https://dlmf.nist.gov/5.15}.

\bibitem{Zeta}
NIST Digital Library of Mathematical Functions,
Section 25.11, \emph{Hurwitz Zeta Function}.
\url{https://dlmf.nist.gov/25.11}.

\bibitem{Lerch}
NIST Digital Library of Mathematical Functions,
Section 25.14, \emph{Lerch's Transcendent}.
\url{https://dlmf.nist.gov/25.14}.

\bibitem{Polylog}
NIST Digital Library of Mathematical Functions,
Section 25.12, \emph{Polylogarithms}.
\url{https://dlmf.nist.gov/25.12}.
\end{thebibliography}
\end{document}
